\subsection{Strongly secant parts and quotients of a Macaulay module}

Let A be a Noetherian ring, M a finitely generated A-module, and S a subset of A. In accordance with the conventions of Chapter VIII, we denote by SM the submodule \(\sum_{s\in S}sM\) of M, and let G be the ideal of A generated by S.

Lemma 2.--- Let \(\overline{\mathfrak{G}}\) be the image of \(\mathfrak{G}\) in A/Ann(M). We have

\[
\operatorname{ht} (\overline {{\mathfrak {G}}}) = \operatorname{codim} (\operatorname{Supp} (\mathrm{M} / \mathrm{SM}), \operatorname{Supp} (\mathrm{M})) .
\]

When moreover SM ≠ M, we have

\[
\operatorname{ht} (\overline {{{\mathfrak {G}}}}) \leqslant \operatorname{Card} (S).
\]

Let a denote the annihilator of M. By the corollary to Proposition 18 of II, § 4, no. 4, the support of the A-module M/SM is V(S + a). Its codimension in Supp(M) is therefore equal to the codimension of V(S + a) in V(a), which in turn equals the codimension of V((S + a)/a) in Spec(A/a), namely the height of G.

Suppose SM ≠ M; the inequality ht(G) ≤ Card(S) is evident when S is infinite, and follows from prop. 4 b) of VIII, § 3, no. 3 when S is finite.

DEFINITION 2.--- Let A be a Noetherian ring, M a finitely generated A-module, and S a finite subset of A. One says that S is strongly secant for M if one has

\[
\operatorname{Card} (\mathrm{S}) \leqslant \operatorname{codim} (\operatorname{Supp} (\mathrm{M} / \mathrm{SM}), \operatorname{Supp} (\mathrm{M})) .
\]

Remarks. - 1) Every finite subset S of A such that SM = M is strongly secant for M. When SM ≠ M, it follows from Lemma 2 that for S to be strongly secant for M, it is necessary and sufficient that Card(S) = codim(Supp(M/SM), Supp(M)), or equivalently ht(G) = Card(S).

\begin{enumerate}
\def\labelenumi{\arabic{enumi})}
\setcounter{enumi}{1}
\tightlist
\item
  If the ring A is local and the module M is nonzero, every subset S of \(m_{A}\) strongly secant for M is secant for M. Indeed, since the A-module M/SM is nonzero, we have
\end{enumerate}

\[
\operatorname{Card} (S) \leqslant \operatorname{codim} (\operatorname{Supp} (M / S M), \operatorname{Supp} (M)) \leqslant \dim (M) - \dim (M / S M)
\]

(VIII, § 1, no. 2, prop. 3 a)), whence our assertion.

PROPOSITION 5. --- Let A be a Noetherian ring, M a finitely generated A-module, and S a finite subset of A. The following conditions are equivalent:

\begin{enumerate}
\def\labelenumi{(\roman{enumi})}
\item
  the subset S of A is strongly secant for M;
\item
  for every element p of Supp(M/SM), the canonical map \(A \to A_{\mathfrak{p}}\) induces a bijection from S onto a subset of \(\mathfrak{p}A_{\mathfrak{p}}\) secant for \(M_{\mathfrak{p}}\) .
\item
  (i) \(\Rightarrow\) (ii): Let \(\mathfrak{p} \in \operatorname{Supp}(\mathrm{M} / \mathrm{SM})\) and let \(S'\) the image of S in \(A_{\mathfrak{p}}\) . The set \(S'\) is contained in the maximal ideal \(\mathfrak{p} A_{\mathfrak{p}}\) and one has
\end{enumerate}

\[
\dim (\mathrm{M} _ {\mathfrak {p}} / \mathrm{S} ^ {\prime} \mathrm{M} _ {\mathfrak {p}}) = \operatorname{codim} (\mathrm{V} (\mathfrak {p}), \operatorname{Supp} (\mathrm{M} / \mathrm{SM}))
\]

(VIII, § 1, no. 4, prop. 9). The inequality Card(S) ≤ codim(Supp(M/SM), Supp(M)) and prop. 3 b) of VIII, § 1, no. 2 imply the relations

\[
\operatorname{Card} (S) + \dim \left(M _ {\mathfrak {p}} / S ^ {\prime} M _ {\mathfrak {p}}\right) \leqslant \operatorname{codim} \left(V (\mathfrak {p}), \operatorname{Supp} (M)\right) = \dim \left(M _ {\mathfrak {p}}\right).
\]

As \(M_{p}\) is nonzero, on the other hand we have \(\dim(M_{\mathfrak{p}}) \leqslant \text{Card}(S') + \dim(M_{\mathfrak{p}}/S'M_{\mathfrak{p}})\) (VIII, § 3, no. 2, formula (8)). Condition (ii) then follows from the inequality \(\text{Card}(S') \leqslant \text{Card}(S)\) .

\begin{enumerate}
\def\labelenumi{(\roman{enumi})}
\setcounter{enumi}{1}
\tightlist
\item
  (ii) \(\Rightarrow\) (i): We may assume SM \(\neq\) M. If condition (ii) is satisfied, then for every prime ideal \(\mathfrak{p}\) of Supp(M/SM)
\end{enumerate}

\[
\operatorname{Card} (S) = \operatorname{Card} \left(S ^ {\prime}\right) \leqslant \dim \left(M _ {\mathfrak {p}}\right) = \operatorname{codim} (V (\mathfrak {p}), \operatorname{Supp} (M)),
\]

which implies (i) by passing to the infimum.

COROLLARY.--- Let A be a Noetherian ring and M a finitely generated A-module. Every M-regular sequence is strongly secant for M.

Let x be an M-regular sequence, and let J be the ideal of \(A\) that it generates. For every prime ideal \(\mathfrak{p} \in \operatorname{Supp}(\mathrm{M/JM})\) , the image of x in \(A_{\mathfrak{p}}\) is an \(\mathrm{M}_{\mathfrak{p}}\)-regular sequence of elements of \(\mathfrak{pA}_{\mathfrak{p}}\) , hence a secant sequence for \(\mathrm{M}_{\mathfrak{p}}\) (VIII, § 3, no. 2, cor. of prop. 3).

PROPOSITION 6. Let A be a Noetherian ring, M a finitely generated Macaulay A-module, and S a finite subset of A strongly secant for M. Then the A-module M/SM is Macaulay.

For every maximal ideal \(\mathfrak{m} \in \operatorname{Supp}(\mathrm{M} / \mathrm{SM})\) , the image of S in \(A_{\mathfrak{m}}\) is secant for \(M_{\mathfrak{m}}\) (prop. 5). Since \(M_{\mathfrak{m}}\) is a Macaulay \(A_{\mathfrak{m}}\)-module, so is \((\mathrm{M} / \mathrm{SM})_{\mathfrak{m}}\) (prop. 4), whence the proposition.

THEOREM 2 (Macaulay-Cohen).--- Let A be a Noetherian ring and M a finitely generated A-module. The following conditions are equivalent:

\begin{enumerate}
\def\labelenumi{(\roman{enumi})}
\item
  the A-module M is Macaulay
\item
  for every ideal J of A generated by an M-regular sequence of elements of A, the A-module M/JM has no embedded associated prime ideals;
\item
  for every finite subset S of A strongly secant for M, the A-module M/SM has no embedded associated prime ideals.
\end{enumerate}

(i)⇒(iii): Let S be a finite subset of A strongly secant for M. The A-module M/SM is Macaulay (prop. 6) and in particular has no embedded associated prime ideals (no. 2, prop. 2, a)).

(iii)⇒(ii): This follows from the cor. of prop. 5.

\begin{enumerate}
\def\labelenumi{(\roman{enumi})}
\setcounter{enumi}{1}
\tightlist
\item
  (ii) → (i): Let \(\mathfrak{p} \in \text{Supp}(M)\) ; let us prove that the \(A_{\mathfrak{p}}\) -module \(M_{\mathfrak{p}}\) is Macaulay. We reason by induction on the integer dim( \(M_{\mathfrak{p}}\) ). If dim( \(M_{\mathfrak{p}}\) ) is zero, \(M_{\mathfrak{p}}\) is an \(A_{\mathfrak{p}}\) -module of finite length, hence Macaulay (example 1, no. 1). Suppose that dim( \(M_{\mathfrak{p}}\) ) is nonzero, that is, that p is not a minimal element of Supp(M) (VIII, § 1, n° 4, prop. 9 a)). Since M has no embedded associated prime ideals, p is contained in no associated prime ideal of M, and there exists an element x of p such that the homothety \(x_{M}\) is injective (§ 1, remark 2). The homothety \(x_{M_{\mathfrak{p}}}\) is then injective, and one has dim( \(M_{\mathfrak{p}}/xM_{\mathfrak{p}}\) ) \textless{} dim( \(M_{\mathfrak{p}}\) ) (VIII, § 3, n° 2, prop. 3). By the induction hypothesis applied to the A-module M/xM and to the prime ideal p of Supp(M/xM), the \(A_{\mathfrak{p}}\) -module \(M_{\mathfrak{p}}/xM_{\mathfrak{p}}\) is Macaulay, which implies that the \(A_{\mathfrak{p}}\) -module \(M_{\mathfrak{p}}\) is Macaulay (n° 3, prop. 4).
\end{enumerate}

